(1) 薄透镜制造公式
\[
\frac1{f(\lambda)}=(n(\lambda)-1)\left(\frac1{R_1}-\frac1{R_2}\right)
=\frac{2[n(\lambda)-1]}R
\]
给
\[
f(\lambda)=\frac{R}{2[n_\infty-1+C/\lambda^2]}.
\]
因 \(n'(\lambda)=-2C/\lambda^3<0\)，所以
\[
f'(\lambda)=-\frac R2\frac{n'(\lambda)}{[n(\lambda)-1]^2}
=\frac{RC}{\lambda^3[n(\lambda)-1]^2}>0.
\]
因此波长越长（红光），折射率越低，焦距越长。

(2) 高斯成像公式为 \(1/s+1/s'=1/f\)。当 \(s>f\) 时成实像，像距与横向放大率为
\[
s'(\lambda)=\frac{s f(\lambda)}{s-f(\lambda)},\qquad
M(\lambda)=-\frac{s'}s=-\frac{f(\lambda)}{s-f(\lambda)}.
\]
将 \(s\) 固定，对波长求导，
\[
\frac{ds'}{d\lambda}=\frac{s^2}{[s-f(\lambda)]^2}f'(\lambda)
=\frac{s^2RC}{\lambda^3[n(\lambda)-1]^2[s-f(\lambda)]^2}>0.
\]
小色散的像距变化受物距与焦距之差的平方放大。

(3) 对 \(\lambda_0\) 附近小谱宽，泰勒展开得到
\[
\Delta s'\simeq
\frac{s^2RC}{\lambda_0^3[n(\lambda_0)-1]^2[s-f(\lambda_0)]^2}\,\Delta\lambda.
\]
若 \(\Delta\lambda>0\)，红光波长更长，焦距与像距也更长，故它的实像焦点离透镜更远。横向放大率 \(M<0\) 表示倒像，且红光放大率绝对值更大。该展开仅在全部所涉及波长仍满足 \(s>f(\lambda)\)、\(\Delta\lambda\) 足够小的范围有效。

(4) 当 \(s\downarrow f(\lambda_0)^+\)，像距 \(s'=sf/(s-f)\to+\infty\)，其对波长的灵敏度也以 \((s-f)^{-2}\) 发散。物体位于精确前焦平面时出射为平行束，不能赋予有限的实像距离；跨过 \(s=f\) 后是虚像分支。若 \(C\to0^+\)，则 \(f(\lambda)\to R/[2(n_\infty-1)]\) 不再随波长变化，纵向色差 \(\Delta s'\to0\)，回到无色散薄透镜。
